\documentclass[11pt]{article}

% ------------------------------------------------------------------
% Speaker notes: explaining the MeanFlow equations on presentation
% slides 30-34 (Stormcast_via_cfm_and_meanflow.pdf). One section per
% slide; each gives the equation followed by suggested bullet points.
% ------------------------------------------------------------------

\usepackage[margin=1in]{geometry}
\usepackage{amsmath, amssymb}
\usepackage{enumitem}
\usepackage{xcolor}
\usepackage{mdframed}
\usepackage[colorlinks=true, linkcolor=blue!50!black, urlcolor=blue!50!black]{hyperref}

% notation shortcuts matching the thesis
\newcommand{\Nbb}{\mathcal{N}}
\newcommand{\Ebb}{\mathbb{E}}

\setlist[itemize]{leftmargin=1.4em, itemsep=2pt, topsep=3pt}

\title{\textbf{Explaining the MeanFlow Equations}\\[2pt]
       \large Speaker bullets for presentation slides 30--34}
\author{}
\date{}

\begin{document}
\maketitle
\vspace{-2.5em}

\noindent The five equations on slides 30--34 build a single argument:
\emph{define} the object we want $\rightarrow$ show \emph{why} it is useful
$\rightarrow$ expose the \emph{catch} $\rightarrow$ make it \emph{computable}
$\rightarrow$ \emph{train} a network on it. The three-panel figure
(Diffusion $\approx 35$ calls / FlowCast $\approx 10$ / MeanFlow $1$) stays
fixed while the equation underneath advances.

% ==================================================================
\section*{Slide 30 --- Average velocity (definition)}
\begin{equation*}
  u(z_r, r, t) \;\triangleq\; \frac{1}{t - r} \int_{r}^{t} v(z_\tau, \tau)\, d\tau
\end{equation*}
\begin{itemize}
  \item \textbf{The picture on the right is the whole story:} diffusion
        curves to the forecast in $\approx 35$ calls, FlowCast straightens it
        to $\approx 10$, MeanFlow wants to jump there in \textbf{1}.
  \item $v(z_\tau, \tau)$ is the \textbf{instantaneous} velocity FlowCast
        already learns --- the local tangent at each point of the path.
  \item $u$ is its \textbf{time-average over the interval $[r, t]$}:
        displacement divided by interval length. It is the single constant
        velocity that carries you from $z_r$ to $z_t$.
  \item Crucially, $u$ is a \textbf{property of the flow itself}, not of any
        network --- and as $t \to r$ it collapses back to $v$
        (one evaluation $=$ FlowCast).
\end{itemize}

% ==================================================================
\section*{Slide 31 --- Single-step transport (why we want $u$)}
\begin{equation*}
  z_t = z_r + (t - r)\, u(z_r, r, t)
\end{equation*}
\begin{itemize}
  \item This is just the definition rearranged --- but read it as a
        \textbf{sampler}: knowing $u$ turns transport into \textbf{one
        algebraic step}, no ODE solver.
  \item Set $(r, t) = (0, 1)$ and a \textbf{single evaluation maps
        noise $\rightarrow$ forecast} (the green dot, one call).
  \item This is where MeanFlow differs from FlowCast in kind, not degree:
        FlowCast \emph{integrates} many small steps; MeanFlow takes
        \textbf{one exact step}.
  \item \textbf{The catch:} $u$ is defined by an intractable path integral
        (slide 30), so we cannot regress a network onto it directly. The cost
        has moved from inference into training.
\end{itemize}

% ==================================================================
\section*{Slide 32 --- The MeanFlow identity \textnormal{(boxed = key result)}}
\begin{equation*}
  u(z_r, r, t) \;=\; v(z_r, r) \;+\; (t - r)\, \frac{d}{dr}\, u(z_r, r, t)
\end{equation*}
\begin{itemize}
  \item We \textbf{kill the integral by differentiating instead of
        integrating.} Multiply slide 30 by $(t - r)$, differentiate in the
        state-time $r$ (product rule $+$ fundamental theorem of calculus).
  \item Result relates the field we \textbf{want} ($u$) to the field we
        \textbf{can sample} ($v$), using only quantities at a \textbf{single
        point of a single trajectory} --- the integral is gone.
  \item The term $(t - r)\, \frac{d}{dr} u$ is the \textbf{correction} that
        lifts the local velocity $v$ to the interval average $u$.
  \item At $t = r$ it reduces to $u = v$ --- \textbf{plain flow matching is
        the boundary condition}, which is exactly why MeanFlow is a clean
        generalisation of FlowCast.
\end{itemize}

% ==================================================================
\section*{Slide 33 --- Expanding the total derivative}
\begin{equation*}
  \frac{d}{dr}\, u(z_r, r, t) \;=\; v(z_r, r)\, \partial_z u \;+\; \partial_r u
\end{equation*}
\begin{itemize}
  \item The $\frac{d}{dr}$ in the identity is a \textbf{total} derivative ---
        the state $z_r$ itself moves along the trajectory as $r$ changes
        ($dz_r / dr = v$).
  \item Chain rule splits it into two \textbf{partials we can compute}:
        motion through space ($v\, \partial_z u$) plus the explicit dependence
        on the state-time ($\partial_r u$). ($t$ is held fixed, so no
        $\partial_t$ term.)
  \item This is the form that maps directly onto a \textbf{Jacobian--vector
        product (JVP)}: one extra forward-mode pass with tangent $(v, 1, 0)$
        for $(z, r, t)$ --- no higher-order gradients.
  \item In short: this line is what makes the identity \textbf{actually
        trainable} on a network.
\end{itemize}

% ==================================================================
\section*{Slide 34 --- Training objective}
\begin{align*}
  \mathcal{L}_{\text{MF}}
  &= \Ebb\Bigl[\, \mathrm{sg}(w)\,
       \bigl\| u_\theta(z_r, r, t, c_t) - \mathrm{sg}(u_{\text{tgt}}) \bigr\|_{2,\beta}^2
     \,\Bigr] + \lambda_{\text{spec}}\, \mathcal{L}_{\text{spec}}, \\[4pt]
  u_{\text{tgt}} &= v + (t - r)\bigl( v\, \partial_z u_\theta + \partial_r u_\theta \bigr)
\end{align*}
\begin{itemize}
  \item We \textbf{regress the network onto the identity's right-hand side}:
        the target $u_{\text{tgt}}$ is just slide 32 with slide 33 substituted
        in, and the network \textbf{bootstraps its own derivative}.
  \item $v = x_1 - x_0$ is the conditional I-CFM velocity (same substitution
        FlowCast makes, valid in expectation) --- \textbf{still teacher-free,
        no diffusion model queried}.
  \item \textbf{stop-gradient ($\mathrm{sg}$)} on the target and on the
        weight: we do not backprop through the JVP, so training stays
        first-order and stable.
  \item \textbf{$w = (\|\Delta\|_{2,\beta}^2 + \varepsilon)^{-p}$} is adaptive
        loss weighting (down-weights hard samples); \textbf{$\|\cdot\|_{2,\beta}$}
        carries the same per-channel weights $\beta = (1, 1, 1, 2)$ and
        \textbf{$\mathcal{L}_{\text{spec}}$} the same \texttt{qpepre} spectral
        term as FlowCast.
  \item A network at zero loss \textbf{satisfies the identity $\rightarrow$
        equals the true average velocity.} And on the $75\%$ of the batch where
        $t = r$, $u_{\text{tgt}} = v$ and this is \textbf{literally the
        FlowCast loss} --- so MeanFlow vs FlowCast is an A/B on the objective
        alone.
\end{itemize}

\vspace{1em}
\begin{mdframed}[backgroundcolor=gray!8, linecolor=gray!40]
\textbf{Presentation note.} The same three-panel figure repeats on all five
slides while the equation underneath advances. If that is intentional
(equations animating in over a fixed visual), these bullets match it. If you
prefer, let the coloured dots / trajectory in the figure light up in step with
the equation --- e.g.\ highlight the single green MeanFlow dot on slide 31 when
you introduce one-step transport.
\end{mdframed}

\end{document}
